\section{Realization in fields and a harvesting model}
\label{sec:sharpness}

The observed controls can be varied independently by local perturbations
through every permitted reference field. We first construct these
perturbations in the original time, then obtain the same rank from natural
parameters in a two-tier harvesting model. A fixed full-rank observation inverse
transfers the configurations proved above to both realizations.

\subsection{Local controls through a specified reference field}
\label{subsec:specified-field}

Fix a nonstationary reference cycle \(\mathcal C\) satisfying the geometric
hypotheses of Section~\ref{sec:passage}, with transported strengths
\(A_{1,0}=-A_{2,0}\ne0\) and actual closing jets
\[
 R_0(0)=0,\qquad R_0'(0)=1,\qquad
 c_0=\tfrac12R_0''(0)\ne0.
\]
Both contacts use positive penetration, and the smooth connectors have
no additional seam event. Choose a second-contact box \(U_2\) and a
smooth closing-arc flow box \(U_3\), with disjoint closures away from
the first contact and return section.
Let \(\mathcal L_{\mathcal C}(U_2,U_3)\) be the finite-dimensional
jointly smooth parameter germs through this field with seam changes
supported in \(\operatorname{int}U_2\), inactive branch fixed there,
and active difference \(h_\lambda B_\lambda\) with smooth
\(B_\lambda\). In \(\operatorname{int}U_3\), which stays away
from the seam, a common smooth field may be added to both branches.
All other data are fixed. The contact controls \(\delta,m\) and
closing controls \(a,b\) have disjoint supports, but their observation
Jacobian retains the couplings below.

\begin{proposition}[Four local controls through every reference field]
\label{prop:specified-rank}
Every such reference field admits a four-parameter unfolding in
\(\mathcal L_{\mathcal C}(U_2,U_3)\) with
\(D_\lambda(a,b,m,\delta)\) invertible at zero. At zero parameters
it recovers the specified field, including its original physical time.
The perturbations preserve matching, embedded regular seams, quadratic
contacts and one common one-turn first-return interval. Every
finite-parameter germ in the allowed class has these common-domain
properties after its own fixed reduction.
\end{proposition}

\begin{proof}
Transport the first inactive integral along the unchanged first
connector. Clock Morse normalization near the second contact gives
\[
 X_0^-=v(t,q)\partial_t,\qquad
 X_0^+=v(t,q)\{[1+(q-t^2)\alpha(t,q)]\partial_t
                    +(q-t^2)\beta(t,q)\partial_q\},\quad v>0.
\]
For a positive smooth \(e\),
\(h_0=e(q-t^2)\) and
\(eB_0=v(\alpha\partial_t+\beta\partial_q)\).
Choose compactly supported cutoffs \(\psi,\zeta\) in \(U_2\),
equal to one on a smaller contact core, and set
\begin{equation}\label{eq:specified-contact-controls}
 h_d=h_0-de\psi,\qquad
 X_{\kappa,d}^+=X_0^++(h_d-h_0)B_0
                  +\kappa h_dv\zeta e^{-1}\partial_q,
 \qquad X_{\kappa,d}^-=X_0^-.
\end{equation}
These changes extend by zero, recover \(X_0^\pm\) at zero, and satisfy
\[
 X_{\kappa,d}^+-X_0^-
 =h_d(B_0+\kappa v\zeta e^{-1}\partial_q).
\]
Thus the traces agree on the moving seam in the original time.
In the core, \(z=q-d\) gives the normalized seam \(z=t^2\) and
coefficients \(v(t,z+d),\alpha(t,z+d),\beta(t,z+d)+\kappa\).
The contact is \((0,d)\), with inactive second seam derivative
\(-2v(0,d)^2<0\) in the normalized defining function, and the
active clock stays positive. Proposition~\ref{prop:quadratic-passage}
therefore gives \(A_2(\kappa,d)=\tfrac43[\beta(0,d)+\kappa]\).
The first passage and connector are unchanged, so
\begin{equation}\label{eq:specified-strength-jacobian}
 \delta=d,\quad m=A_{1,0}+\tfrac43[\beta(0,d)+\kappa],\qquad
 D_{(\kappa,d)}(m,\delta)(0)=
 \begin{pmatrix}4/3&(4/3)\beta_q(0,0)\\0&1\end{pmatrix}.
\end{equation}
The second strength is already transported; its incoming threshold is
\(\Gamma_1^{-1}(d)\).

In \(U_3\), use actual base flow time \(s\) and the second out-port
coordinate transported by that flow as \(z\). Then
\(X_0^-=\partial_s\), the entrance holonomy is \(q\mapsto z=q\),
and the remaining transition is \(R_0\). Add to both branches
\[
 Y_\eta=\chi(s)\xi(z)(\eta_0+\eta_1z)\partial_z,
 \qquad \chi\ge0,\qquad \int\chi(s)\,ds=1,
\]
with compact support and \(\xi=1\) on a fixed transverse interval.
Small inputs and controls stay in that interval, so the exact box
holonomy and closing transition are
\[
 H_\eta(z)=e^{\eta_1}z+\eta_0\phi(\eta_1),\qquad
 \phi(w)=\int_0^1e^{wt}\,dt,\qquad R_\eta=R_0\circ H_\eta.
\]
Consequently
\begin{equation}\label{eq:specified-four-jacobian}
 D_{(\eta_0,\eta_1,\kappa,d)}(a,b,m,\delta)(0)=
 \begin{pmatrix}
 1&0&0&0\\2c_0&1&0&0\\
 0&0&4/3&(4/3)\beta_q(0,0)\\0&0&0&1
 \end{pmatrix},\qquad \det=4/3.
\end{equation}
Disjoint supports make the closing leg independent of \((\kappa,d)\).
The inverse function theorem gives smooth original-field coordinates
\((a,b,m,\delta)\), with the higher profiles still coupled.

The seam $q-t^2-d\psi(t,q)=0$ has $q$ derivative greater than $1/2$
for small $d$ and agrees with the old seam on a boundary collar. The
reference has a positive angular clock on a fixed tube; matching gives
common local Lipschitz bounds for the perturbed physical fields.
Choose cores where the cutoffs are one, ports with strict transverse margins,
and strictly inactive connecting arcs, including the cutoff collars.
Proposition~\ref{prop:quadratic-passage} gives the first entry and moving
first exit in each core, with second penetration $\Gamma_1(q)-d$.
Lemma~\ref{lem:common-tube} therefore gives the actual one-turn return
\[
 P_\lambda(q)=R_\eta\{d+
       \Gamma_{2,\kappa,d}(\Gamma_1(q)-d)\}.
\]
For any finite-parameter germ in the allowed class, matching, smooth
parameter dependence and the contact implicit function theorem preserve
these same strict clock, core and port margins after a germ-dependent
reduction. The common-return conclusion needs no observation rank.
\end{proof}

\begin{corollary}[Cyclicity in the allowed local perturbation class]
\label{cor:specified-cyclicity}
The cyclicity of every reference cycle of
Proposition~\ref{prop:specified-rank}, relative to
\(\mathcal L_{\mathcal C}(U_2,U_3)\), is exactly six: it is the
supremum, over finite-parameter germs in that class, of the numbers
of distinct isolated periodic orbits approaching \(\mathcal C\)
at parameters tending to zero.
Its constrained cyclicities on $m=0$, on $\delta=0$, and on their
intersection are respectively four, four and two.
\end{corollary}

\begin{proof}
Each germ has the common actual geometry after its own reduction,
so Proposition~\ref{prop:common-count}, restricted to its observations,
gives at most six fixed points.
Every periodic orbit approaching \(\mathcal C\) meets that section
interval. An increasing return has no nontrivial finite-period orbit:
if \(P(q)>q\) or \(P(q)<q\), its successive iterates are strictly
ordered. Thus fixed points count these cycles. Conversely,
Proposition~\ref{prop:specified-rank} and
Proposition~\ref{prop:single-scale-six} give six simple cycles
approaching \(\mathcal C\).
Restricting the same observation inverse to the three indicated slices,
Proposition~\ref{prop:sharp-slice-attainment} attains their bounds from
Corollary~\ref{cor:four-slices}.
\end{proof}
The reduction is per finite-parameter germ. A prescribed model or
fixed-seam subfamily can lack the seam-height direction and requires
its own rank calculation.

\subsection{A two-tier escapement-harvesting model}
\label{sec:harvesting}

The continuous policy of Bohn, Rebaza and Speer
\cite[Model~Two, equation~(17)]{bohn2011} stops prey harvest below
a lower threshold, increases it linearly between two thresholds,
and caps it above the upper one. Onana Ndzana et al.
\cite[equations~(5)--(6)]{onana2019} use it in a delayed
Leslie--Gower model with prey refuge. We intensify harvest above
the upper threshold, using the normalized smooth
predator-harvesting core of Huang, Gong and Ruan
\cite[equation~(9) and Theorem~3.6]{huang2013}:
\begin{equation}\label{eq:harvesting:model}
 \dot x=x(1-x-y)-u(x-L)_+-v(x-U)_+,\qquad
 \dot y=y\left(d-\beta\frac yx\right)-h.
\end{equation}
Here \(x,y>0\) are prey and predator densities,
\(\beta,d,h,u,v>0\), and \(0<L<U<1\).
Here \(h\) denotes predator harvest; the multiplier defect in the
preceding analysis is \(-b\). The predator parameters \(d,\beta,h\) describe growth, density
limitation relative to prey abundance, and constant-yield harvest.
The marginal prey efforts are \(0,u,u+v\). We work in a compact
subset of the positive quadrant. The middle smooth extension is
\begin{equation}\label{eq:harvesting:middle}
 F_c=x(1-x-y)-u(x-L),\qquad
 G=y(d-\beta y/x)-h,\qquad X_c=(F_c,G),
\end{equation}
and the exterior extensions are
\[
 X_\ell=(x(1-x-y),G),\qquad X_r=(F_c-v(x-U),G).
\]
They agree on the whole seams \(x=L,U\). The middle extension is
inactive at both contacts, with active sides outside \(L\le x\le U\).
The branches are analytic for \(x>0\) and the physical field is
locally Lipschitz.

\begin{theorem}[Natural harvesting realization]
\label{thm:harvesting-realization}
There exists \(\rho_0>0\) such that, for every fixed
\(0<\rho<\rho_0\), there exists \(u_0(\rho)>0\) with the following
property. For every fixed \(0<u<u_0(\rho)\), positive parameters
\(\beta,d,h,v\) and \(0<L<U<1\) give an embedded neutral reference
cycle \(\mathcal C\) of \eqref{eq:harvesting:model}, with exactly two
quadratic contacts at its minimum and maximum prey densities. In one
transported ordinary return coordinate the contact levels agree and
\[
 A_1=A_0=-A_2\ne0,\qquad
 R(0)=0,\qquad R'(0)=1,\qquad c_0=\tfrac12R''(0)<0.
\]
The observation map \((a,b,m,\delta)\) has rank four in the natural
family \((\beta,d,h,u,v,L,U)\). The local cyclicity is exactly six,
with constrained cyclicities four, four and two on \(m=0\), on
\(\delta=0\), and on their intersection. For each fixed \(\rho,u\),
the fixed full-rank observation inverse also gives the main-scale
saturation and actual quartic conclusions of
Corollary~\ref{cor:realization-transfer}.
All constants and scale reductions may depend on that reference.
\end{theorem}

We construct a smooth cyclic fold, fit the thresholds to its prey
extrema, and cancel the transported strengths with a positive effort
ratio. The directions \((\nu,L,v,U)\), with \(\beta,u\) fixed,
then give the observation inverse. The scale \(\varepsilon\) is
introduced only after fixing this reference and differs from the
effort \(u\).

\subsection{Threshold fitting, cancellation and closing curvature}

Appendix~\ref{app:harvesting-fold} constructs actual cyclic folds of
the smooth core, for every sufficiently small fixed radius \(\rho>0\).
The fold is embedded and nonstationary, lies in the positive quadrant,
and has exactly one nondegenerate prey minimum and maximum. In radial
coordinates its displacement satisfies
\eqref{eq:harvesting:fold-data}, with \(\ell<0\), and the
\((\beta,\nu)\) Jacobian of \((D,D_r)\) is invertible.
These strict properties allow the thresholds and positive efforts to be
fitted to a fixed fold.

For positive \(u\), use
\begin{equation}\label{eq:harvesting:perturbed-chart}
 y_e=\tfrac45-u(1-5L),\qquad
 d=8\beta+\tfrac15+\nu,\qquad h=y_e(d-5\beta y_e).
\end{equation}
The middle field vanishes at \((1/5,y_e)\). At \(u=0\) the chart
reduces to \eqref{eq:harvesting:core-chart} and is independent of
\(L\). Fix a positive fold radius in
\eqref{eq:harvesting:fold-data}. The invertible \((\beta,\nu)\)
Jacobian of \((D,D_r)\) continues a fold \(\mathcal C(u,L)\)
for small \(u,L-L_0\), where \(L_0\) is the unperturbed prey
minimum. Its curvature and two nondegenerate extrema persist.
For its minimum \(\lambda_-(u,L)\),
\(\partial_L\lambda_-(0,L_0)=0\), so
\begin{equation}\label{eq:harvesting:threshold-fit}
 L=\lambda_-(u,L),\qquad U=\max_x\mathcal C(u,L(u))
\end{equation}
has a unique local smooth solution with \(0<L(u)<U(u)<1\).
After reducing \(u\), all parameters in
\eqref{eq:harvesting:perturbed-chart} remain positive. The fitted
orbit stays in the closed middle strip, touching its boundary only
at the extrema, and is a physical orbit for every \(v\). There
\(F_c=0\), \(F_{c,y}=-x\), and nonstationarity gives \(G\ne0\),
so
\begin{equation}\label{eq:harvesting:quadratic-contact}
 \ddot x=F_{c,y}G=-xG\ne0.
\end{equation}

Figure~\ref{fig:harvesting-orbit} illustrates this fitted threshold geometry.
\begin{figure}[tb]
 \centering
 \includegraphics[width=90mm]{figures/harvesting-orbit.pdf}
 \caption{Computed geometry of \eqref{eq:harvesting:model} for
 \(u=10^{-6}\), \(\beta\simeq0.031443254\),
 \(d\simeq0.451546877\), and \(h\simeq0.260619080\).
 The floating-point middle orbit touches the fitted vertical
 thresholds \(L\simeq0.191154841\) and \(U\simeq0.210000000\)
 at its prey extrema. Open circles mark the contacts \(p_1,p_2\),
 arrows show increasing physical time, and the dot marks the middle
 equilibrium. This computed picture illustrates the fitted geometry;
 it does not certify the neutral reference or its six-cycle unfolding.}
 \label{fig:harvesting-orbit}
\end{figure}

Positive penetration is outward, into \(x<L\) or \(x>U\).
Transport one inactive first integral \(I\) from the first box through
the connector to the second, with outward direction positive. These are
local first-integral charts joined along the selected connecting tube;
no single-valued first integral around the full periodic orbit is assumed. At
\(p_i=(\xi_i,y_i)\), put \(\sigma_1=-1\), \(\sigma_2=1\),
\(s_i=I_x(p_i)\). Then \(I_y(p_i)=0\) and
\(\sgn s_i=\sigma_i\). Clock Morse normalization of
\(\mathfrak h_i=\sigma_i(x-\xi_i)\), preserving \(I\), gives
\begin{equation}\label{eq:harvesting:clock}
 \mathfrak h_i=e_i(t,I)(I-t^2),\qquad
 e_i(0,0)=\frac{\sigma_i}{s_i}>0,\qquad
 \dot t(0,0)=\omega_i
 =\sqrt{-s_iF_{c,y}(p_i)G(p_i)/2}>0.
\end{equation}
These clocks differ from the core frequency \(\omega\). For active
difference \(\mathfrak h_iB_i\), the normalized transverse coefficient is
\[
 \beta_{i,\mathrm{pass}}(0,0)
 =\frac{e_i(0,0)}{\omega_i}dI(B_i)
 =\frac{\sigma_i(B_i)_x}{\omega_i}.
\]
Here \(B_1=-u\partial_x\), \(B_2=-v\partial_x\); thus
\eqref{eq:passage:strength-jet} gives
\begin{equation}\label{eq:harvesting:strengths}
 A_1=K_1u=\frac{4u}{3\omega_1}>0,\qquad
 A_2=-K_2v=-\frac{4v}{3\omega_2}<0.
\end{equation}
The opposing passage signs come from the opposite outward directions at
the two thresholds; both harvesting efforts are positive.
The inactive clocks and transport are independent of \(v\), so
the exact positive ratio
\begin{equation}\label{eq:harvesting:effort-relation}
 v=u\frac{K_1}{K_2}=u\frac{\omega_2}{\omega_1}
\end{equation}
cancels the strengths in the same transported coordinate.
Both contacts lie on one inactive orbit, so their levels agree.
Their passage derivatives equal one at the touch and the middle fold
has multiplier one. Lemma~\ref{lem:passage:finite-inverse} gives
\(Q(q)=q+q_+^{5/2}\mathcal L(\sqrt{q_+})\), with zero ordinary
quadratic defect. The regular map \(R\) is the inactive monodromy.

Its curvature is negative in the outward-positive coordinate. First
set \(u=0\) and fix the cyclic-fold parameters at a small radius
\(\rho\). Vary the initial radial coordinate \(r\) with these
parameters fixed. Smooth finite-time normal-form flow has prey component
\[
 x(r,\theta)=\tfrac15+rp(\theta)+O(r^2+\rho^2r),
\]
where \(p\) is a nonzero first harmonic and the \(r\)-derivative
of the remainder is \(O(r+\rho^2)\). Its unique nearby prey minimum
\(M(r)\) satisfies \(M'(\rho)=-\gamma+O(\rho)<0\), \(\gamma>0\).
The minimum locus \(F_c=0\) is transverse, since its derivative
along the orbit is \(\ddot x\ne0\). If \(H(r)\) is the
transported inactive integral there, \(I_y=0\), \(I_x<0\) give
\(H'(\rho)=I_xM'(\rho)>0\). For the radial core return \(P_c\)
at these fixed parameters,
\(R=H\circ P_c\circ H^{-1}\). At the neutral fixed point,
\[
 R''(0)=\frac{P_c''(\rho)}{H'(\rho)}<0
\]
by \eqref{eq:harvesting:fold-data} and \(\ell<0\).
Continuation to fitted small positive \(u\) preserves this strict
sign and \(H'>0\). The cancelling composition has zero quadratic
defect, so \(c_0<0\) for the physical reference as well.

\subsection{Natural rank and the actual first return}

Use \((\nu,L,v,U)\) in \eqref{eq:harvesting:perturbed-chart}, with
\(\beta,u\) fixed. Choose the first incoming coordinate from the
middle extension and lower seam only. The regular monodromy \(R\)
is independent of \(v,U\): changes of the upper ports alter the
intermediate inactive holonomies but preserve their middle-return
composition. The level difference \(\delta\) is independent of
\(v\). At the upper tangency \(F_c(U,y_2(U))=0\) and \(I_y=0\),
so
\begin{equation}\label{eq:harvesting:upper-rank}
 \delta_U=I_x(p_2)\ne0,\qquad m_v=-K_2\ne0.
\end{equation}
Hence
\begin{equation}\label{eq:harvesting:rank-matrix}
 D_{(\nu,L,v,U)}(a,b,m,\delta)=
 \begin{pmatrix}
 a_\nu&a_L&0&0\\ b_\nu&b_L&0&0\\
 *&*&-K_2&*\\ *&*&0&\delta_U
 \end{pmatrix}.
\end{equation}
For the upper-left minor first set \(u=0\) at the fixed core fold.
The lower tangency is \(p_1(L)=(L,1-L)\); its level in an unchanged
core transverse coordinate has derivative
\begin{equation}\label{eq:harvesting:lower-level}
 \zeta_L=dI(p_1'(L))=I_x(p_1)\ne0.
\end{equation}
Choose this coordinate to agree with the incoming one at the base,
and center the contact by \(T_L(\zeta(L))=0\). Then
\(R_L=T_L\circ P_c\circ T_L^{-1}\). Since at the fold
\(P_c(\zeta)=\zeta\) and \(P_c'(\zeta)=1\),
\begin{equation}\label{eq:harvesting:lower-rank}
 a_L=0,\qquad b_L=P_c''(\zeta)\zeta_L
                   =2c_0\zeta_L\ne0.
\end{equation}
Coordinate-motion terms cancel at the fixed point and unit multiplier;
another transverse scaling supplies only a nonzero factor. Contact
motion leaves the offset derivative in \(\nu\) unchanged because
\(D_r=0\), and \eqref{eq:harvesting:scaled-return} gives
\(a_\nu\ne0\). The minor is therefore
\(a_\nu\,2c_0\zeta_L\ne0\). These objects extend smoothly to
\(u=0\), so at each fixed radius the minor persists for small
positive \(u\), with fitted thresholds and
\eqref{eq:harvesting:effort-relation}. All couplings remain in
\eqref{eq:harvesting:rank-matrix}, whose determinant is the minor
times \((-K_2)\delta_U\ne0\). The chart directions embed in
\((\beta,d,h,u,v,L,U)\), proving natural rank four.

The fitted cycle is compact, embedded and nonstationary. On a fixed tube,
the vertical seams stay regular and the matched model has common local
Lipschitz bounds. The positive angular clock, transverse ports and strictly
inactive connectors persist under small parameter changes. In each core,
\eqref{eq:harvesting:clock} and the smooth active difference satisfy the
actual passage hypotheses, including first entry and first moving exit.
Lemma~\ref{lem:common-tube} gives the chronological first return on one
common interval. This verifies the geometric assumptions independently
of the rank computation.

\subsection{Transfer to physical parameters}

\begin{corollary}[One observation inverse for the core results]
\label{cor:realization-transfer}
Fix a rank-four reference constructed above and fix its remaining physical
coordinates. Use the fixed full-rank observation inverse and the common
physical return interval \(I\). The sharp bounds \(6,4,4,2\) of
Theorem~\ref{thm:six-main} hold in that family. In the sector
$A_1(0)=B>0$, $c_0=-C<0$, including every harvesting reference,
Theorem~\ref{thm:main-saturation-summary} gives weighted necessity,
the \(\mathcal K_4\) cluster restriction, realization of every compact
subset of \(\operatorname{int}\mathcal K_4\), and Hausdorff convergence.
Theorem~\ref{thm:actual-quartic} and Proposition~\ref{quartic:strata}
give the actual quadruple cycle and its local quartic strata, with ordinary
multiplicities. Off the multiple-root set, the total count in \(I\)
is two, four or six, and each count occurs on a nonempty open set.
Exactly two simple negative-state companions persist throughout this
quartic neighbourhood; at the quadruple point there are three distinct cycles.
All harvesting parameters remain positive and $0<L<U<1$; constants and
scale thresholds may depend on the fixed ecological reference.
\end{corollary}

\begin{proof}
The inverse function theorem applied to the Jacobians
\eqref{eq:specified-four-jacobian} and \eqref{eq:harvesting:rank-matrix}
gives a smooth observation inverse; Lemma~\ref{lem:common-tube}
gives the common return interval. In the harvesting family the inverse maps
\((a,b,m,\delta)\) to \((\nu,L,v,U)\), with \(\beta,u\) fixed.
Fix this full-rank observation inverse before introducing the scale.
Every compact scaled control box enters its domain for sufficiently small
scale, while positivity and threshold ordering persist.
Theorem~\ref{thm:six-main} gives the sharp counts, including their
attainment on the three slices. Theorem~\ref{thm:main-saturation-summary}
applies to this fixed inverse. Theorem~\ref{thm:actual-quartic} localizes
all roots in \(I\) and prepares the positive cluster; its two negative
companions and Proposition~\ref{quartic:strata} give the stated local counts
and ordinary multiplicities.
Since the increasing return has no nontrivial finite-period orbit, these
are precisely the local cycle counts. This completes
Theorem~\ref{thm:harvesting-realization}.
\end{proof}
